\documentclass{article}
\usepackage[T1]{fontenc}
\usepackage{lmodern}
\usepackage{graphicx}
\usepackage{amsmath,amssymb}
\usepackage[a4paper,margin=2cm]{geometry}
\usepackage[absolute,overlay]{textpos}
\setlength{\TPHorizModule}{1mm}
\setlength{\TPVertModule}{1mm}
\usepackage[indonesian]{babel}
\usepackage{hyperref}
\setlength{\emergencystretch}{2em}
% unit: Halaman 8

\begin{document}

% === Halaman 8 (llm) ===

\begin{itemize}
    \item Jika plainteks/cipherteks hanya menggunakan 26 huruf alfabet (A--Z), maka entropi Shannon maksimum terjadi ketika seluruh 26 huruf memiliki probabilitas yang sama:
    \[ p_i = \frac{1}{26} \]
    sehingga
    \begin{align*}
    H_{max} &= - \sum_{i=1}^{26} p_i \log_2 p_i \\
    &= - \sum_{i=1}^{26} \frac{1}{26} \log_2 \frac{1}{26} = \log_2 26 = 4.7004 \text{ bit/karakter}
    \end{align*}
    Jadi, untuk teks yang hanya terdiri dari 26 huruf alfabet, entropi Shannon maksimum adalah sekitar 4,7004 bit/karakter.
    
    \item Jika pesan terdiri dari 256 karakter ASCII, maka entropi maksimum adalah $\log_2 256 = 8$
    
    \item Makin besar entropi, makin sulit memecahkan cipherteks.
\end{itemize}

\end{document}
